\documentclass[../../../include/open-logic-section]{subfiles}

\begin{document}

\olfileid{sth}{choice}{tarskiscott}
\olsection{তার্স্কি--স্কট কৌশল}

\olref[cardinals][cardsasords]{defcardinalasordinal}-এ
অঙ্কবাচক সংখ্যাকে ক্রমসংখ্যা হিসেবে সংজ্ঞায়িত করেছি।
সেজন্য সুক্রমবিন্যাসের স্বতঃসিদ্ধ ধরেছিলাম। একমাত্র
কারণ, এটি এই বিষয়ে প্রচলিত মানক রীতি।

তাই সুক্রমবিন্যাস নিয়ে দার্শনিক প্রশ্নে যাওয়ার আগে
বোঝা দরকার: আমরা প্রচলিত রীতি থেকে \emph{সরে এসে},
সুক্রমবিন্যাস না ধরেও অঙ্কবাচক সংখ্যার তত্ত্ব গড়তে
\emph{পারি}। \olref[sfr][siz][]{chap}-এর
$\cardeq{A}{B}$, $\cardle{A}{B}$ এবং
$\cardless{A}{B}$-র সংজ্ঞা তখনও ব্যবহার করা যায়।
শুধু \emph{অঙ্কবাচক সংখ্যার} নতুন ধারণা লাগবে।

সরল মনে হতে পারে, $A$-র অঙ্কবাচকতা এভাবে সংজ্ঞায়িত করি:
\[
	\Setabs{x}{\cardeq{A}{x}}.
\]
\olref[cardinals][hp]{sec}-এর একেবারে শেষে আলোচিত
ফ্রেগের $\# x Fx$-সংজ্ঞার সঙ্গে এটি তুলনা করতে পারো।
কিন্তু সেখানে ইঙ্গিত করা কারণেই এই সংজ্ঞা চলে না।
যেকোনো একপদী সেট $\{\emptyset\}$-র সমসংখ্যক।
অথচ স্তরক্রমের প্রত্যেক উত্তরসূরি পর্যায়ে নতুন একপদী
সেট তৈরি হয়—আগের পর্যায়টির একপদী সেটই ধরো।
তাই $\Setabs{x}{\cardeq{A}{x}}$-র কোনো স্তরাঙ্ক
থাকতে পারে না; এটি সেট নয়।

এই অসুবিধা কাটাতে তার্স্কি ও স্কটের একটি কৌশল
ব্যবহার করি:\footnote{মনে রেখো: আলাদা করে নিষেধ
না থাকলে সব সূত্রে পরামিতি থাকতে পারে।}

\begin{defn}[তার্স্কি--স্কট]
যেকোনো সূত্র $\phi(x)$-র জন্য $[ x : \phi(x)] $
হলো $\phi(x)$ মেনে চলে এমন সম্ভাব্য ক্ষুদ্রতম স্তরাঙ্কের
সব $x$-এর সেট (আর কোনো $\phi$-সন্তুষ্টকারী না থাকলে
$\emptyset$)।
\end{defn}

সংজ্ঞাটি সার্থক কি না দেখা দরকার। $\ZF$-এ কাজ করি।
\olref[spine][foundation]{zfentailsregularity}
অনুযায়ী প্রত্যেক $x$-এর $\setrank{x}$ আছে।
এখন কোনো বস্তু যদি $\phi$ মেনে চলে, তবে এমন
ক্ষুদ্রতম স্তরাঙ্ক $\alpha$ নিই যার জন্য
$(\exists x\subseteq V_\alpha)\phi(x)$; অর্থাৎ
$(\forall \beta\in \alpha)(\forall x \subseteq V_\beta)
\lnot \phi(x)$। তখন পৃথকীকরণ দিয়ে
$[x : \phi(x)]$-কে
$\Setabs{x \in V_{\alpha+1}}{\phi(x)}$
হিসেবে সংজ্ঞায়িত করা যায়।

তার্স্কি--স্কট কৌশলের ন্যায্যতা পাওয়ায় এবার এর
সাহায্যে অঙ্কবাচকতা সংজ্ঞায়িত করতে পারি:

\begin{defn}
$A$-র \textsc{ts}-অঙ্কবাচকতা হলো
$\text{tsc}(A) = [x : \cardeq{A}{x}]$।
\end{defn}

\textsc{ts}-অঙ্কবাচক সংখ্যার সংজ্ঞায় সুক্রমবিন্যাস
লাগে না। তবু ওই স্বতঃসিদ্ধ ছাড়াই দেখানো যায় যে
\emph{\textsc{ts}-অঙ্কবাচক সংখ্যাগুলি}
\olref[cardinals][cardsasords]{defcardinalasordinal}-এ
সংজ্ঞায়িত \emph{অঙ্কবাচক সংখ্যার} মতো অনেকটাই
আচরণ করে। যেমন,
\olref[cardinals][cardsasords]{lem:CardinalsBehaveRight}
ও \olref[card-arithmetic][opps]{lem:SizePowerset2Exp}
\textsc{ts}-অঙ্কবাচক সংখ্যার ভাষায় বললে তাদের
প্রমাণ $\ZF$-এই চলে; সুক্রমবিন্যাস ধরতে হয় না।

এই প্রসঙ্গে আরও লক্ষণীয়, তার্স্কি--স্কট কৌশলে
ক্রমসংখ্যারও একটি তত্ত্ব গড়া যায়।
$\tuple{A, <}$ সুক্রমবিন্যাস হলে
$\text{tso}(A, <) = [\tuple{X,
R} : \ordeq{\tuple{A, <}}{\tuple{X, R}}]$
ধরি। অঙ্কবাচক ও ক্রমসংখ্যার এই রীতি সম্বন্ধে আরও
জানতে দেখো \citet[chs.~9--12]{Potter2004}।

\end{document}
